If a molecule is composed of identical atoms\footnote{The two atoms must belong not merely to the same element, but also to the same isotope of that element.}, its Hamiltonian is also invariant under interchange of the coordinates of the two nuclei. A term is called symmetric with respect to the nuclei if its wave function is unchanged when the nuclei are interchanged, and antisymmetric if the wave function changes sign. Symmetry with respect to the nuclei is closely connected with the parity and the sign of the term. Interchanging the nuclear coordinates is equivalent to reversing the coordinates of all particles (electrons and nuclei), followed by reversal of the electronic coordinates alone. It follows that a term is symmetric with respect to the nuclei if it is even (odd) and at the same time positive (negative). If, on the other hand, the term is even (odd) and simultaneously negative (positive), it is antisymmetric with respect to the nuclei.
